\chapter{Tutorial 2: Shell Commands, Streams, grep and Vim}
\begin{center}
\textit{January 10, 2018}
\end{center}

\section{List of Commands}

\begin{center}
\begin{tabular}{@{}lp{0.65\textwidth}@{}}
\toprule
Command & What it does \\
\midrule
\code{cd} & change directory; with no argument goes home; \code{cd -} goes back
            to the previous directory \\
\code{clear} & clear the screen \\
\code{ls} & list files, e.g.\ \code{ls ..}, \code{ls -l}, \code{ls -a} \\
\code{touch} & create an empty file (or update a file's timestamp) \\
\code{pwd} & display the present working directory \\
\code{cat} & print the contents of a file \\
\code{head} & print the first 10 lines (\code{head -n 5} for the first 5) \\
\code{tail} & print the last 10 lines \\
\code{sort} & sort lines in lexicographical order \\
\code{uniq} & remove \emph{adjacent} duplicate lines (so usually used after \code{sort}) \\
\code{shuf} & output a random permutation of the input lines \\
\code{tac} & like \code{cat} but prints the lines in reverse order \\
\code{wc} & count lines, words and characters \\
\bottomrule
\end{tabular}
\end{center}

\section{Input and Output Streams}
Think of a program as a black box with one input stream and two output streams:
\begin{itemize}
  \item \textbf{stdin}: by default the keyboard.
  \item \textbf{stdout}: by default the screen. stdout is \emph{buffered}: output
        is collected and written in blocks, since writing character by character
        is expensive.
  \item \textbf{stderr}: by default also the screen. stderr is \emph{not}
        buffered, so error messages appear immediately. Example: the nuclear
        reactor is about to explode --- we want to see that message \emph{now},
        not when the buffer fills up.
\end{itemize}
We will solidify these concepts when we do I/O in C++.

\begin{lstlisting}[style=bash]
g++ printer.cc -o print
./print > out.txt 2> err.txt   # stdout to out.txt, stderr to err.txt
./print &> all.txt             # all streams to all.txt
./print > all.txt 2>&1         # same: send stderr where stdout is going
\end{lstlisting}

\begin{important}[Stream numbers]
stdin is 0, stdout is 1, stderr is 2. \code{2>\&1} redirects stderr to the same
place as stdout; \code{\&>} redirects both. Order matters: \code{2>\&1} must come
\emph{after} \code{> file}.
\end{important}

\paragraph{Pipes.} \code{cmd1 | cmd2} connects stdout of \code{cmd1} to stdin of
\code{cmd2}.

\paragraph{Embedded commands.} \code{\$(cmd)} is replaced by the output of
\code{cmd}, e.g.\ \code{egrep "\$(cat pattern)" words.txt}.

\section{grep and Regular Expressions}
\code{grep} (and \code{egrep}, which uses extended regular expressions) prints
the lines of its input that match a pattern:
\begin{lstlisting}[style=bash]
egrep "int" wordcollection.txt
\end{lstlisting}

\begin{center}
\begin{tabular}{@{}cl@{}}
\toprule
Symbol & Meaning \\
\midrule
\code{\^{}} & match the beginning of the line \\
\code{\$} & match the end of the line \\
\code{.} & match any single character \\
\code{?} & the preceding item 0 or 1 times \\
\code{*} & the preceding item 0 or more times \\
\code{+} & the preceding item 1 or more times \\
\code{[ ]} & any one of the characters inside \\
\bottomrule
\end{tabular}
\end{center}

\section{Vim Commands}
\begin{center}
\begin{tabular}{@{}ll@{}}
\toprule
Keys & Effect \\
\midrule
\code{i} & enter insert mode (\code{Esc} returns to normal mode) \\
\code{Esc :w} & write (save) \\
\code{Esc :wq} & write and quit \\
\code{Esc :x} & save (only if there are changes) and quit \\
\code{Esc :xa} & do \code{:x} on all open buffers/tabs \\
\code{Esc :q!} & quit without saving \\
\code{u} & undo \\
\code{Ctrl-r} & redo \\
\code{dd}, \code{yy}, \code{p} & delete line, copy (yank) line, paste \\
\code{/text} & search forward for \code{text} \\
\bottomrule
\end{tabular}
\end{center}

\section*{Summary}
Know the basic file/text commands, the three standard streams and how to
redirect them (including \code{2>\&1}), pipes and \code{\$(...)}, basic regex
for \code{egrep}, and enough Vim to edit, save and quit.
